% capitolo 5: Risultati dei modelli GARCH e delle correlazioni
\chapter{Risultati dei modelli GARCH e delle correlazioni}
\label{cap:risultati}

\section{Selezione del modello e stima dei parametri}
L'analisi empirica ha inizio con la selezione del modello ottimale per analizzare le dinamiche della varianza condizionata, per poi passare alla stima e interpretazione 
dei parametri nel modello selezionato. La valutazione si basa sulla minimizzazione del Criterio di Informazione Bayesiano (BIC), confrontando le strutture sGARCH, 
GJR-GARCH ed EGARCH, stimate con distribuzione Gaussiana e Student-t. Come mostrato nella tabella \ref{tab:bic_valori}, il modello asimmetrico EGARCH-t minimizza il valore 
per tutti e 6 gli asset in esame.

\begin{table}[htbp]
    \centering
    \caption{Valori del BIC}
    \label{tab:bic_valori}
    \resizebox{\textwidth}{!}{%
    \begin{tabular}{lcccccc}
        \toprule
        \textbf{Asset} & \textbf{sGARCH\_norm} & \textbf{GJR-GARCH\_norm} & \textbf{EGARCH\_norm} & \textbf{sGARCH\_std} & \textbf{GJR-GARCH\_std} & \textbf{EGARCH\_std} \\
        \midrule
        Oro & 2.6817 & 2.6722 & 2.6677 & 2.6190 & 2.6152 & 2.6086 \\
        Nucor & 4.3930 & 4.3903 & 4.3759 & 4.3039 & 4.3041 & 4.2960 \\
        GM & 4.3607 & 4.3561 & 4.3307 & 4.2492 & 4.2453 & 4.2368 \\
        Apple & 3.8663 & 3.8441 & 3.8228 & 3.7610 & 3.7500 & 3.7397 \\
        S\&P 500 (SPY) & 2.6115 & 2.5867 & 2.5623 & 2.5374 & 2.5121 & 2.5007 \\
        USD/CNH & 0.2254 & 0.2214 & 0.2113 & 0.1146 & 0.1157 & 0.1091 \\
        \bottomrule
    \end{tabular}%
    }
\end{table}

Questa superiorità viene confermata anche dal Criterio di Informazione di Akaike (AIC) e dalle tabelle delta (dBIC e dAIC), i cui risultati completi sono riportati in 
Appendice. Impostando l'EGARCH-t come benchmark ($\Delta=0$), tutti gli altri modelli restituiscono delta ampiamente positivi e mostrano come la Student-t migliori 
l'adeguatezza in tutti i modelli. Definita la specifica ottimale, la Tabella \ref{tab:parametri_egarch} riporta le stime dei parametri del modello EGARCH-t.

\begin{table}[htbp]
    \centering
    \caption{Stima dei parametri del modello EGARCH-t}
    \label{tab:parametri_egarch}
    \resizebox{\textwidth}{!}{%
    \begin{tabular}{lcccccc}
    \toprule
    \textbf{Parametro} & \textbf{Oro} & \textbf{Nucor} & \textbf{GM} & \textbf{Apple} & \textbf{S\&P 500 (SPY)} & \textbf{USD/CNH} \\
    \midrule
    $\mu$ & 0.059*** & 0.068* & 0.058** & 0.112*** & 0.070*** & -0.001 \\
    $\omega$ & -0.002 & 0.014*** & 0.013*** & 0.031*** & -0.018*** & -0.057*** \\
    $\theta^*$ (segno) & 0.036*** & -0.023** & -0.043*** & -0.092*** & -0.165*** & 0.019 \\
    $\alpha^*$ (ampiezza) & 0.102*** & 0.088*** & 0.101*** & 0.175*** & 0.191*** & 0.200*** \\
    $\beta^*$ & 0.988*** & 0.990*** & 0.992*** & 0.968*** & 0.970*** & 0.979*** \\
    $\nu$ & 5.353*** & 5.261*** & 4.157*** & 4.622*** & 6.099*** & 4.445*** \\
    \bottomrule
    \end{tabular}%
    }
    \vspace{0.15cm}
    \footnotesize{\textit{Nota:} *, **, *** indicano significatività al 10\%, 5\% e 1\% (basata su errori standard classici). Le stime si riferiscono al modello includente le variabili dummy, i cui coefficienti sono omessi per brevità; il $\beta^*$ riportato è quello impiegato per i calcoli nella Sezione 5.2. Rispetto a \textit{rugarch}: $\theta^*$ = alpha1, $\alpha^*$ = gamma1, $\omega$ = costante.}
\end{table}

Il parametro di asimmetria ($\theta^*$) evidenzia il classico effetto leva per i titoli azionari (GM e Apple) e per l'indice di mercato S&P 500, mentre per Nucor risulta 
significativo solo al 10\%. I valori inferiori a 0 indicano che gli shock negativi incrementano la volatilità in misura maggiore rispetto a shock positivi di pari entità. 
Per quantificare la dimensione di questa asimmetria, si osservi l'indice S&P 500: l'impatto dei rialzi pesa solamente 0,026 ($\alpha^* + \theta^*$), mentre quello dei ribassi 
pesa ben 0,356 ($\alpha^* - \theta^*$). Al contrario, l'Oro presenta un coefficiente positivo (0,036), evidenziando un'asimmetria inversa, diversa da quella azionaria. 
La valutazione della sua effettiva natura di bene rifugio viene pertanto demandata all'analisi delle correlazioni nei periodi di stress, mentre per il tasso di cambio USD/CNH 
non si rifiuta la simmetria. 

I gradi di libertà della distribuzione ($\nu$) si attestano tra 4 e 6, confermando la presenza di code grasse (\emph{fat tails}) e l'inadeguatezza di una distribuzione 
normale. Infine, il coefficiente di persistenza ($\beta^*$) risulta prossimo a 1 per tutti gli asset (superiore a 0.96). Calcolando l'\gls{emivita} (\emph{half-life}) degli 
shock tramite la formula $\log(0.5)/\log(\beta^*)$, emerge che i movimenti si dimezzano in tempi variabili: da un assorbimento più rapido di circa 21 giorni per Apple, 
a un rientro lento di quasi 86 giorni per General Motors, evidenziando un'elevata persistenza della volatilità.